\SourceLeaf{237}%
the first method of describing the curved surface. Keeping the notation of Article 4, we also set:
\[
\begin{array}{lll}
\dfrac{\mathrm{d}\mathrm{d}W}{\mathrm{d}x^2}=P',&\dfrac{\mathrm{d}\mathrm{d}W}{\mathrm{d}y^2}=Q',&\dfrac{\mathrm{d}\mathrm{d}W}{\mathrm{d}z^2}=R'\\[6pt]
\dfrac{\mathrm{d}\mathrm{d}W}{\mathrm{d}y\mathord{.}\mathrm{d}z}=P'',&\dfrac{\mathrm{d}\mathrm{d}W}{\mathrm{d}x\mathord{.}\mathrm{d}z}=Q'',&\dfrac{\mathrm{d}\mathrm{d}W}{\mathrm{d}x\mathord{.}\mathrm{d}y}=R''
\end{array}
\]
so that
\[
\begin{aligned}
\mathrm{d}P&=P'\mathrm{d}x+R''\mathrm{d}y+Q''\mathrm{d}z\\
\mathrm{d}Q&=R''\mathrm{d}x+Q'\mathrm{d}y+P''\mathrm{d}z\\
\mathrm{d}R&=Q''\mathrm{d}x+P''\mathrm{d}y+R'\mathrm{d}z
\end{aligned}
\]
Now, since \(t=-\frac{P}{R}\), we find by differentiation that
\[
\begin{aligned}
RR\mathrm{d}t&=-R\mathrm{d}P+P\mathrm{d}R=(PQ''-RP')\mathrm{d}x\\
&\quad +(PP''-RR'')\mathrm{d}y+(PR'-RQ'')\mathrm{d}z
\end{aligned}
\]
or, on eliminating \(\mathrm{d}z\) by means of the equation \(P\mathrm{d}x+Q\mathrm{d}y+R\mathrm{d}z=0\),
\[
\begin{aligned}
R^3\mathrm{d}t&=(-RRP'+2PRQ''-PPR')\mathrm{d}x\\
&\quad +(PRP''+QRQ''-PQR'-RRR'')\mathrm{d}y
\end{aligned}
\]
In exactly the same way, we obtain
\[
\begin{aligned}
R^3\mathrm{d}u&=(PRP''+QRQ''-PQR'-RRR'')\mathrm{d}x\\
&\quad +(-RRQ'+2QRP''-QQR')\mathrm{d}y
\end{aligned}
\]
From this we conclude that
\[
\begin{aligned}
R^3T&=-RRP'+2PRQ''-PPR'\\
R^3U&=PRP''+QRQ''-PQR'-RRR''\\
R^3V&=-RRQ'+2QRP''-QQR'
\end{aligned}
\]
Substituting these values in the formula of Article 7, we obtain the following symmetric expression for the measure of curvature \(k\):
\[
\begin{aligned}
&(PP+QQ+RR)^2k\\
&\quad=PP(Q'R'-P''P'')+QQ(P'R'-Q''Q'')+RR(P'Q'-R''R'')\\
&\qquad+2QR(Q''R''-P'P'')+2PR(P''R''-Q'Q'')+2PQ(P''Q''-R'R'')
\end{aligned}
\]
\Article{10}
We obtain a still more complicated formula, made up of fifteen elements, if we follow the second general method for describing the nature of
